%Dokument pro XeLaTeX, překlad pomocí XeLaTeX
\documentclass[a4paper,oneside,11pt,usenames,dvipsnames]{article}

%XeLaTeXové věci a fonty
\usepackage{fontspec,xltxtra,xunicode}
\defaultfontfeatures{Scale=MatchLowercase}

%Písma
%\setromanfont[Mapping=tex-text]{Linux Libertine O}
%\setsansfont[Mapping=tex-text]{Linux Biolinum O}
%\setmonofont[Mapping=tex-text]{Linux Libertine Mono O}

%Pro správné dělení českých slov
\usepackage{polyglossia}
\setdefaultlanguage{czech} %TODO: Případně změňte jazyk
%České uvozovky buď přímo nebo jako ,,uvozovky``

%Některé speciální znaky
\usepackage{textcomp}
% Url
\usepackage{url}
\def\UrlFont{\ttfamily\slshape}

%Barvy pro hyperref
\usepackage{color}
% Hyperref pro odkazy v PDF a URL
\usepackage{hyperref}
% Hyperref a PDF nastavení
\hypersetup{
    pdftitle={Titulek},    % title TODO: Vyplnit titulek
    pdfauthor={Nějaký Autor},     % author TODO: Vyplnit autora
    colorlinks=true,       % false: boxed links; true: colored links
    linkcolor=RoyalBlue,          % color of internal links
    citecolor=RoyalBlue,        % color of links to bibliography
    filecolor=RoyalBlue,      % color of file links
    urlcolor=RoyalBlue           % color of external links
}

%Nadpisy, záhlaví a zápatí
\usepackage[pagestyles]{titlesec}

%Nadpisy bezpatkovým písmem
\titleformat{\part}[display]{\filright}{\Large\sffamily Část\ \thepart}{0.25em}{\bfseries\sffamily\huge\upshape}
\titleformat{\section}[block]{\sffamily\Large\bfseries}{\thesection}{1em}{\Large}
\titleformat{\subsection}[block]{\sffamily\large\bfseries}{\thesubsection}{1em}{\large}
\titleformat{\subsubsection}[block]{\sffamily\normalsize\bfseries}{\thesubsubsection}{1em}{\normalsize}

\begin{document}

\thispagestyle{empty}
\begin {center}
	\vspace *{10mm}
	{\Huge \sffamily \bfseries Výstižný titulek práce}

	\vspace *{20mm}
	{\Large \sffamily \bfseries Nějaký autor}

	\vspace *{7.5mm}
	{\large \sffamily Nějaký datum}
\end {center}

\begin{abstract}
	Abstrakt Text informativního abstraktu max. 200 slov, np.:
	Technologie X je to a ono, používá se tam a tam. Experiment takový a takový přinesl tyto výsledky.
\end{abstract}

\setlength{\parskip}{1ex plus 0.25ex minus 0.125ex}
  
\section{Úvod}
\label{sec:uvod}

Text úvodu. Mělo by tu být zhruba:
\begin{itemize}
	\item Stručná charakteristika použitých AI technik.
	\item Stručná charakteristika řešené/modelované domény.
	\item Cílem této sekce je vysvětlit základy toho, s čím pracujete, a nasměrovat čtenáře na vhodnou literaturu s detaily.
\end{itemize}

Celkový rozsah tohoto dokumentu je cca 1200 slov (bez příloh).

Pro PDF výstup: \verb|xelatex soubor.tex|.

Struktura dokumentu je podle formátu \href{https://projects.ncsu.edu/eslglobe/nmswishe/IMRAD\%20Handout.pdf}{IMRAD}.

Je-li sekce dále členěná, měl by tu být ,,obsahový odstavec``: V sekci \ref{subsec:uvod-o-tom} popíšu to, v sekci \ref{subsec:uvod-o-onom} zas ono.

\subsection{Sekce o tom}
\label{subsec:uvod-o-tom}

Text.

Odkazy na zdroje se po vyplnění bibliografické citace v sekci Reference vkládají do textu přes \verb|\cite{xxx}|. ,,Nějaká zajímavá myšlenka z jiné knihy`` \cite{autor-jina-kniha}. Většinou je ale práce s literaturou o vybrání důležitých myšlenek a jejich stručném shrnutí vlastními slovy. Parafráze zajímavých myšlenek z nějaké knihy, viz \cite{autor-nejaka-kniha}.

\subsection{Sekce o onom}
\label{subsec:uvod-o-onom}

Text.
 
\section{Metody}
\label{sec:metody}

Text. Mělo by tu být zhruba:
\begin{itemize}
	\item Popis postupu (návrh a realizace/experimenty), jak jste se dostali k výsledku z následující sekce.
	\item Případné využití generativní AI by také mělo být zdokumentováno zde, detaily promptů v přílohách.
	\item Cílem této sekce je umožnit čtenáři, aby provedl to, co vy a mohl si ověřit, že dojde ke stejnému/srovnatelnému výsledku.
\end{itemize}

\section{Výsledky}
\label{sec:vysledky}

Text. Mělo by tu být zhruba:
\begin{itemize}
	\item Popis výsledku (vytvořený model, chování systému, konkrétní měření experimentů), tj. co vám vyšlo, když jste provedli postup z předchozí sekce.
	\item Je vhodné odkázat na detaily v příloze, ale neměl by tu být jen odkaz.
	\item Cílem sekce je komunikovat čtenáři výsledek práce a umožnit mu srovnání při replikaci.
\end{itemize}

Obrázky, tabulky a podobně se vkládají nejprve jako plovoucí objekt. V textu by měla být zmínka o obrázku s křížovým odkazem na jeho značku, np.: Na obrázku \ref{fig:obrazek-neceho} je schéma něčeho. Křížové odkazy lze nastavit, aby zobrazovaly i stránku, což se může u plovoucích objektů hodit.

\begin{figure}
	\centering
	%\includegraphics[width=1.0\linewidth]{cesta-k-obrazku.jpg}
	Sem se teprve vloží vlastní obrázek pomocí \verb|includegraphics|. Obvykle je potřeba nastavit nějaké rozumné rozměry, nejlépe relativně k šířce řádku a zarovnat obrázek na střed.
	\caption{Stručný popisek obrázku – pokud je přejatý pak s odkazem \cite{malir-obrazek-neceho}.\label{fig:obrazek-neceho}}
\end{figure}


\section{Diskuze}
\label{sec:diskuze}

Text. Mělo by tu být zhruba:
\begin{itemize}
	\item Komentář získaných výsledků (co znamenají?)
	\item Zasazení výsledků do širšího kontextu (odpovídá to dalším znalostem z literatury?)
	\item Přiznání možných limitů (co by šlo udělat lépe?)
	\item Možnosti dalšího rozšíření.
	\item Cílem sekce je výsledek zhodnotit a představit v kontextu.
\end{itemize}

\section{Závěr}
\label{sec:zaver}

Text závěru. Poslední sekcí také může být ,,Diskuze a závěr``.

\begin{thebibliography}{11}
\setlength{\parskip}{0ex}

	\bibitem{autor-nejaka-kniha} \textsc{Autor, N.}: Nějaká kniha.

	\bibitem{autor-jina-kniha} \textsc{Autor, J.}: Jiná kniha.

	\bibitem{autor-elektronicky} \textsc{Autor, D.}: Nějaký elektronický zdroj. \textless{}\url{https://www.nekde.xx/neco}\textgreater{}.

	\bibitem{malir-obrazek-neceho} \textsc{Malíř, N.}: Obrázek něčeho.

	\bibitem{selinux-by-example} \textsc{Mayer, F. -- MacMillan, L. -- Caplan, K.:} \textit{SELinux by Example: Using Security Enhanced Linux.} 1$^{st}$\,Ed. Upper Saddle River, Prentice Hall 2006. 456\,s. ISBN-10: 0-13-196369-4.

	\bibitem{wiki-access-control} \textsc{Wikipedia:} \textit{Access Control} [online]. Poslední změna 2009-03-24 [cit. 2009-03-24]. Dostupný z~WWW: \textless\url{http://en.wikipedia.org/wiki/Access_Control}\textgreater.

\end{thebibliography}

\appendix
\section*{Přílohy}

Text příloh – typicky rozsáhlejší tabulky, obrázky, zdrojové kódy, (lze i formou externích souborů), vždy okomentované a odkázané z hlavního textu.

\end{document}
